
\section{Related Work}

\subsection{Low-Rank Matrix Factorization}

\textbf{Classical SVD.} The Eckart--Young theorem (Eckart \& Young, 1936) states 
that the best rank-$r$ approximation of $W$ in Frobenius norm is the truncated 
SVD $W_r = U_{:, :r} \Sigma_{:r, :r} V_{:r, :}^\top$. Adaptive SVD builds 
on this but selects $r$ via an energy threshold rather than a validation 
loop or fixed value.

\textbf{Rank selection heuristics.} Classical rank selection includes AIC 
(Akaike, 1974), BIC (Schwarz, 1978), and cross-validation. For SVD-based 
compression specifically, prior work uses reconstruction error thresholds 
(A3), fixed budget ratios (SDRM), or optimal hard thresholding on 
singular values under noise (Gavish \& Donoho, 2014). Our criterion is a 
\textbf{cumulative energy threshold}, distinct from per-singular-value thresholds.

\textbf{Randomized SVD.} Halko et al. (2011) provide a randomized algorithm with 
$\mathcal{O}(d r \log k)$ complexity. For our layer sizes ($\le$3072$$\times$768), exact 
SVD is faster and simpler.

\subsection{Activation-Aware Compression}

\textbf{AWQ.} Lin et al. (2023) observe that activation magnitude correlates with 
weight importance and use it to scale quantization. Applied to low-rank 
factorization, this becomes the A3 method (Fisher diagonal weighting).

\textbf{SmoothQuant.} Xiao et al. (2023) migrate activation scale to weights for 
quantization. Analogous to AASVD's covariance weighting for low-rank.

\textbf{A3.} Uses Fisher diagonal $F_i = \mathbb{E}[x_i^2]$ to weight the SVD input 
space. Requires calibration. Adaptive SVD is calibration-free.

\textbf{AASVD.} Uses full covariance $C_X = \mathbb{E}[x x^\top]$. Requires 
streaming EMA accumulation. Adaptive SVD is calibration-free.

\textbf{Key distinction.} Adaptive SVD does \textbf{not} use activation statistics. Its 
rank criterion is intrinsic to the weight matrix.

\subsection{Parameter-Efficient Fine-Tuning}

\textbf{LoRA.} Hu et al. (2021) add trainable rank-$r_L$ updates $B A$ with $B = 0$. 
Adaptive SVD uses the same residual structure, integrated with a frozen SVD base.

\textbf{AdaLoRA.} Zhang et al. (2023) adaptively allocate LoRA rank per layer based 
on singular value importance. Our adaptation is on the \textbf{base factorization}, 
not the adapter.

\textbf{PiSSA.} Meng et al. (2024) initialize LoRA adapters with principal singular 
components. Adaptive SVD uses the top-$r$ components as the \textbf{frozen base} 
and zero-initialized LoRA as the residual. These are complementary strategies.

\textbf{OFT.} Qiu et al. (2023) use orthogonal fine-tuning. Different primitive.

\subsection{Attention-Specific Compression}

\textbf{Head pruning.} Voita et al. (2019) prune entire heads. Adaptive SVD preserves 
head count and compresses per-head projections via low-rank factorization.

\textbf{GQA.} Ainslie et al. (2023) reduce KV head count during pretraining. Adaptive 
SVD preserves head count but compresses per-head projection matrices.

\textbf{KV cache compression.} H2O (Zhang et al., 2023), StreamingLLM (Xiao et al., 
2024) reduce cache size at inference. Different target.

\subsection{Position of Our Work}

\begin{table}[H]
\centering
\small
\begin{tabular}{llllll}
\toprule
\textbf{Method} & \textbf{Calib} & \textbf{Weighted} & \textbf{Rank} & \textbf{Per-layer rank} & \textbf{T5 bias} \\
\midrule
SVD fixed rank \& No \& No \& Fixed \& No \& No \\
SDRM \& No \& No \& Budget \& No \& No \\
A3 \& Yes \& Fisher diag \& Escalation \& Partial \& No \\
AASVD \& Yes \& Full cov \& Fixed \& No \& No \\
\textbf{Adaptive SVD} & \textbf{No} & \textbf{No} & \textbf{Energy} & \textbf{Yes} & \textbf{Yes} \\
\bottomrule
\end{tabular}
\end{table}

\textbf{Key distinction.} Adaptive SVD is the only method that:
\begin{enumerate}
\item Requires no calibration data.
\item Selects rank per-projection via an energy criterion.
\item Preserves T5 relative attention bias.

\end{enumerate}
\subsection{Comparison Table}
